% TOP_PROBLEM: {"id": 5300065, "title": "Measure and dimension of transcendental parameter hairs", "category_id": 11, "category": {"id": 11, "name": "dynamical_systems", "display_name": "Dynamical Systems", "description": "Problems about long-term behavior of deterministic systems, Hamiltonian dynamics, and periodic orbits.", "slug": "dynamical-systems", "order_index": 11, "created_at": "2026-07-31T15:26:25.671Z"}, "status": "partially_solved"}
% FIRST_POSTED: null
% ENTRY_KIND: statement_only
% Imported from: batch06.tex; UnsolvedMath ID 5300065
\documentclass[11pt]{article}
\usepackage[margin=1in]{geometry}
\usepackage{amsmath,amssymb,mathtools}
\usepackage{fontspec,unicode-math}
\setmainfont{Latin Modern Roman}
\setsansfont{Latin Modern Sans}
\setmathfont{Latin Modern Math}
\usepackage{xurl,hyperref}
\hypersetup{colorlinks=true,linkcolor=blue,urlcolor=blue}
\newcommand{\sref}[2]{\hyperlink{source:#1:#2}{[S#2]}}
\newcommand{\cataloguescope}[1]{\paragraph{Mathematical question.}#1}
\begin{document}
% UnsolvedMath ID 5300065; source catalogue code AMR-052-0065
\setcounter{section}{5300064}
\section{Measure and dimension of transcendental parameter hairs}\label{5300065}
{\small\sffamily UnsolvedMath ID 5300065 · Dynamical Systems}
\subsection{Definitions and mathematical statement}

\paragraph{Shared notation.}$\mathbb N=\{1,2,\ldots\}$, and $\mathbb Z,\mathbb Q,\mathbb R,\mathbb C$ denote the integers, rationals, reals and complex numbers. Any other conventions are specified locally.


An entire function is a holomorphic map $f:\mathbb C\to\mathbb C$. Its iterates are $f^0=\mathrm{id}$ and $f^{n+1}=f\circ f^n$. Its Fatou set is the open set where the family of iterates is normal (in the spherical metric), and its Julia set $J(f)$ is the complement. The exponential and trigonometric parameter families here are
\[E_\lambda(z)=\lambda e^z,\quad S_\lambda(z)=\lambda\sin z,\quad C_\lambda(z)=\lambda\cos z,\qquad \lambda\in\mathbb C\setminus\{0\}.\]Devaney describes parameter hairs as the Cantor families of curves in the complex parameter plane along which the corresponding map has Julia set equal to $\mathbb C$. Measure means planar Lebesgue measure, and dimension means Hausdorff dimension. These are hairs in parameter space, distinct from Cantor bouquets in the dynamical $z$-plane.
\paragraph{Statement recorded in the supplied source.}
Determine the measure and Hausdorff dimension of the parameter hairs for each of $E_\lambda$, $S_\lambda$ and $C_\lambda$. \sref{5300065}{2}
\subsection{Short English statement}
Determine how large the parameter hairs are, in measure and Hausdorff dimension, for all three exponential, sine and cosine families.
\subsection{Sources}
\begin{description}
\item[\hypertarget{source:5300065:1}{S1}] ULAM.AI and the UnsolvedMath Contributors, UnsolvedMath: A Curated Collection of Open Mathematics Problems, record 5300065 (AMR-052-0065). CC BY 4.0. \url{https://huggingface.co/datasets/ulamai/UnsolvedMath}.
\item[\hypertarget{source:5300065:2}{S2}] Ben Bielefeld, Mikhail Lyubich and contributors, \emph{Problems in Holomorphic Dynamics}, arXiv:math/9205209v1 (1992), Robert Devaney, ``Open Questions in Non-Rational Complex Dynamics'', \S4, seventh Problem, printed p. 36. \url{https://arxiv.org/abs/math/9205209v1}.
\end{description}
\label{end:5300065}
\end{document}
